\chapter{Headroom Context Compression Pipeline}
\label{chap:compression}

\section{The Context Window Bottleneck}
Long-running conversational interactions and multi-step tool executions rapidly accumulate tokens. Context window bloat results in two severe degradations:
\begin{enumerate}
    \item \textbf{Inference Latency}: Prompt evaluation time ($O(N)$ to $O(N^2)$ in attention layers) grows linearly or quadratically, degrading Time-to-First-Token (TTFT).
    \item \textbf{Loss of Salience}: Models experience "needle-in-a-haystack" degradation as context becomes congested with verbose tool outputs, repeated boilerplate, redundant whitespace, and syntax structures.
\end{enumerate}

To eliminate this bottleneck, Finch.ai integrates \texttt{headroom-ai}, a context-compression engine that intercepts prompt messages before they reach the model.

\section{Compression Strategy and Guardrails}
The compression pipeline is managed by \texttt{HeadroomPipeline} in \texttt{ai\_assistant/compression/pipeline.py}. It applies semantic-preserving transformations to text, tool payloads, JSON structures, and programming code.

\subsection{System Prompt Protection}
The system prompt defines Finch.ai's personality, operational constraints, and tool usage protocols. Any semantic mutation of system instructions risks jailbreaking or drift. Finch.ai strictly guarantees:
\begin{equation}
    \texttt{compress\_system\_messages} = \text{False}
\end{equation}
The system message is bypassed in the pipeline and forwarded byte-for-byte to the inference engine.

\subsection{Recent Turn Preservation}
The user's most recent query and the immediately preceding context contain the highest informational recency. Finch.ai configures:
\begin{equation}
    \texttt{compression\_protect\_recent} = 2
\end{equation}
This ensures that the last two dialogue turns retain full lexical fidelity, while older conversational history and previous tool outputs are aggressively compressed.

\subsection{Activation Threshold}
Short conversational exchanges do not require compression overhead. The configuration defines:
\begin{equation}
    \texttt{compression\_min\_tokens} = 100
\end{equation}
When the aggregate token count is below this threshold, compression is bypassed, preserving zero-overhead execution for brief queries.

\section{Applied Transformations}
The pipeline dynamically evaluates content type and applies tailored compression strategies:
\begin{itemize}
    \item \textbf{Code AST \& Whitespace Truncation}: Strips redundant comments, collapses repeated indentation, and compresses import statements while retaining function signatures and logic.
    \item \textbf{JSON \& Structured Output Compaction}: Converts verbose, formatted JSON into dense key-value representations, discarding duplicate schemas.
    \item \textbf{Log \& Terminal Output Compaction}: Collapses repeated stack traces, progress bars, and status updates, keeping only error messages and concluding status lines.
    \item \textbf{Natural Language Pruning}: Eliminates conversational filler, redundant acknowledgments, and repetitive turns.
\end{itemize}

\section{Real-Time Telemetry and Token Economics}
During each inference pass, the pipeline records token statistics. Let $T_{\text{orig}}$ denote the original token count and $T_{\text{comp}}$ denote the compressed token count. The metrics reported are:
\begin{align}
    \Delta T &= T_{\text{orig}} - T_{\text{comp}} \\
    \text{Savings \%} &= \left( \frac{T_{\text{orig}} - T_{\text{comp}}}{T_{\text{orig}}} \right) \times 100\%
\end{align}

Table~\ref{tab:compression_performance} presents empirical compression results across distinct message types captured in Finch.ai benchmarks.

\begin{table}[htbp]
\centering
\small
\begin{tabular}{lcccc}
\toprule
\textbf{Payload Type} & \textbf{Original Tokens} & \textbf{Compressed Tokens} & \textbf{Tokens Saved} & \textbf{Reduction} \\
\midrule
Terminal Build Log & 1,420 & 412 & 1,008 & 70.99\% \\
JSON API Response & 890 & 325 & 565 & 63.48\% \\
Python AST / Code Snippet & 680 & 390 & 290 & 42.65\% \\
Multi-Turn Chat History & 1,150 & 780 & 370 & 32.17\% \\
Short User Prompt & 45 & 45 & 0 & 0.00\% (Protected) \\
\bottomrule
\end{tabular}
\caption{Empirical Headroom Context Compression Benchmarks}
\label{tab:compression_performance}
\end{table}

\section{Offline Fallback Mechanisms}
Token counting relies primarily on \texttt{tiktoken} (\texttt{o200k\_base} or \texttt{cl100k\_base}). In offline environments or air-gapped systems where tokenizer vocabulary downloads are unavailable, Finch.ai implements an automated fallback to heuristic token estimation ($\approx 4$ characters per token for English text and code), preventing runtime exceptions.
